\section{Introduction}
\label{sec:intro}

In our previous report, we demonstrated that value-based methods, such as Rainbow DQN, can successfully solve pixel-based continuous control tasks. However, adapting Q-learning to continuous domains requires action space discretization and even complex branching architectures to scale to more complex tasks. A more mathematically natural approach lies in policy gradient and actor-critic methodologies, which directly output continuous action distributions and are the standard choice for robotic locomotion. 

In this study, we evaluate two state-of-the-art continuous control algorithms: Proximal Policy Optimization (PPO) \cite{schulman2017ppo} and Soft Actor-Critic (SAC) \cite{haarnoja2018soft}. While these algorithms represent the theoretical gold standard for physical control tasks, applying them directly to high-dimensional visual observations introduces deep practical challenges. Specifically, we found both algorithms to be highly sensitive to their underlying hyperparameter configurations. Even on the relatively simple Hopper environment, establishing a stable learning signal proved remarkably difficult. For instance, PPO exhibits extreme sensitivity to its target Kullback-Leibler (KL) divergence penalty; a threshold set too high causes the value loss to explosively diverge, while a threshold set too low results in premature policy collapse and a complete failure to learn.

To navigate this extreme hyperparameter sensitivity, we pivot our methodology from manual heuristic tuning to automated, large-scale search using the Optuna optimization framework~\cite{akiba2019optuna}. Given the computationally expensive nature of pixel-based rendering, committing to long training horizons for unverified configurations is intractable. Instead, we designed an optimization pipeline that executes dozens of short-horizon trials. By evaluating the early learning dynamics over fewer environment steps, Optuna allows us to rapidly prune configurations that catastrophically diverge or stagnate. This broad but shallow search strategy effectively isolates the narrow, stable hyperparameter regimes required before scaling up to full training horizons.